% ---------- ¶12  contribution 2 + scope disclaimer ----------
Our paper also provides an information-based account of the residential energy-efficiency gap
\citep{allcott_is_2012, gerarden_assessing_2017}. The standard diagnosis appeals to split
incentives, credit constraints, and inattention to future energy costs \citep{allcott_consumers_2011, sallee_rational_2014, allcott_evaluating_2015}. We identify a
complementary channel internal to the disclosure game. A seller holding a poor rating can
renovate and publish a good one, earning the quality premium, or stay silent and earn the
non-disclosure premium instead, so the premium to opacity looks like an opportunity cost of
renovation that raises the private hurdle above engineering cost. In equilibrium it does not, except in the corner where every informed low-quality seller hides. This has an unexpected policy
implication. Because the equilibrium absorbs enforcement through the fraction of sellers who
conceal, the renovation threshold depends only on the price gap between disclosed quality
levels and is independent of enforcement intensity. Renovation subsidies, enforcement against
concealment, and investment in audit coverage act on three distinct margins and do not substitute for one another,
contrary to the conventional treatment in
\citet{gillingham_palmer_2014}. The result holds outside the corner in which every informed low-quality seller hides, and with one concealment cost per market. % [AUTHOR REWRITE] The extension of Section~\ref{ss:attention} brings the inattention
diagnosis inside the disclosure game: buyers who attend less to energy when no rating is shown let
silence earn a premium that no belief about quality can produce. We do not claim to have estimated the one true model of housing
disclosure; we propose one way of reading these data through a workhorse framework, and we are
explicit throughout about which parameters are point-identified and which are bounded.
